\chapter*{نمایه‌ی موضوعی}
\addcontentsline{toc}{chapter}{نمایه‌ی موضوعی}
\markboth{نمایه‌ی موضوعی}{نمایه‌ی موضوعی}

شماره‌صفحه‌های این نمایه به محل معرفی یا بررسی اصلی هر موضوع اشاره
می‌کنند. عنوان‌ها به همان بخش در متن پیوند داده شده‌اند؛ بنابراین در
نسخه‌ی دیجیتال می‌توان مستقیماً به محل مربوط رفت.

\section*{مفاهیم و اصطلاح‌ها}
\begin{description}
  \item[آنتروپی شانون] \hyperref[sub:shannon-entropy]{از غافلگیری تا آنتروپی شانون}، ص.\ \pageref{sub:shannon-entropy}
  \item[اطلاعات در طبیعت] \hyperref[sec:information-in-nature]{اطلاعات در طبیعت}، ص.\ \pageref{sec:information-in-nature}
  \item[اطلاعات و معنا] \hyperref[sub:meaning-versus-information]{معنای پیام با مقدار اطلاعات یکی نیست}، ص.\ \pageref{sub:meaning-versus-information}
  \item[اطلاعات متقابل] \hyperref[sub:entropy-and-message-information]{آنتروپی با اطلاعات یک پیام چه تفاوتی دارد؟}، ص.\ \pageref{sub:entropy-and-message-information}
  \item[بار الکتریکی] \hyperref[sub:charge-and-electric-field]{بار الکتریکی و میدان الکتریکی}، ص.\ \pageref{sub:charge-and-electric-field}
  \item[بازنمایی ذهنی] \hyperref[sec:mental-representation]{بازنمایی ذهنی}، ص.\ \pageref{sec:mental-representation}
  \item[بازنمایی لایه‌ای] \hyperref[sec:layered-representation]{ماهیت لایه‌ای بازنمایی}، ص.\ \pageref{sec:layered-representation}
  \item[بیت] \hyperref[sub:logic-levels]{آیا صفر و یک واقعاً دو ولتاژ مشخص‌اند؟}، ص.\ \pageref{sub:logic-levels}
  \item[پیچیدگی کولموگروف] \hyperref[sub:kolmogorov-complexity]{آنتروپی شانون و پیچیدگی کولموگروف}، ص.\ \pageref{sub:kolmogorov-complexity}
  \item[داده] \hyperref[sec:what-no-definition-says]{آنچه هیچ تعریفی مستقیماً نمی‌گوید}، ص.\ \pageref{sec:what-no-definition-says}
  \item[داده، اطلاعات، دانش و خرد] \hyperref[sec:levels-of-data]{هرم DIKW}، ص.\ \pageref{sec:levels-of-data}
  \item[خازن و سلف] \hyperref[sub:capacitors-and-inductors]{خازن و سلف}، ص.\ \pageref{sub:capacitors-and-inductors}
  \item[ظرفیت کانال] \hyperref[sub:channel-capacity]{چرا نمی‌توان هرقدر خواست اطلاعات فرستاد؟}، ص.\ \pageref{sub:channel-capacity}
  \item[شیطان ماکسول] \hyperref[sub:maxwells-demon]{شیطان ماکسول}، ص.\ \pageref{sub:maxwells-demon}
  \item[سیگنال] \hyperref[sub:what-is-signal]{سیگنال چیست؟}، ص.\ \pageref{sub:what-is-signal}
  \item[سیگنال آنالوگ] \hyperref[sub:analog-signal]{سیگنال آنالوگ چیست؟}، ص.\ \pageref{sub:analog-signal}
  \item[سیگنال دیجیتال] \hyperref[sub:digital-signal]{سیگنال دیجیتال: گسسته بودن به چه معناست؟}، ص.\ \pageref{sub:digital-signal}
  \item[نمونه‌برداری] \hyperref[sub:sampling]{نمونه‌برداری: ثبت لحظه‌های جداگانه}، ص.\ \pageref{sub:sampling}
  \item[نویز] \hyperref[sub:what-is-noise]{نویز دقیقاً چیست؟}، ص.\ \pageref{sub:what-is-noise}
  \item[ولتاژ] \hyperref[sub:voltage]{ولتاژ: اختلاف انرژی به ازای هر واحد بار}، ص.\ \pageref{sub:voltage}
  \item[قانون اهم] \hyperref[sub:resistance-circuit-ohms-law]{مقاومت، مدار و قانون اهم}، ص.\ \pageref{sub:resistance-circuit-ohms-law}
  \item[اصل لاندائر] \hyperref[sub:landauer-principle]{اصل لاندائر: پاک کردن اطلاعات چه هزینه‌ای دارد؟}، ص.\ \pageref{sub:landauer-principle}
\end{description}

\section*{اشخاص و نظریه‌ها}
\begin{description}
  \item[ارسطو] \hyperref[sub:aristotle-model]{ارسطو: مدل داخلی مفید}، ص.\ \pageref{sub:aristotle-model}
  \item[افلاطون] \hyperref[sub:plato-cave]{افلاطون: غار و بازنمایی‌های ناقص}، ص.\ \pageref{sub:plato-cave}
  \item[دکارت] \hyperref[sub:descartes-doubt]{شک در بازنمایی‌های حسی}، ص.\ \pageref{sub:descartes-doubt}
  \item[کانت] \hyperref[sub:kant-framework]{ذهن چارچوب تحمیل می‌کند}، ص.\ \pageref{sub:kant-framework}
  \item[پیرس] \hyperref[sub:peirce-triad]{سه‌گانه‌ی پیرس: شمایل، نمایه و نماد}، ص.\ \pageref{sub:peirce-triad}
  \item[شرودینگر] \hyperref[sec:information-in-nature]{اطلاعات در طبیعت}، ص.\ \pageref{sec:information-in-nature}
  \item[شانون] \hyperref[sub:communication-model]{مدل ارتباط}، ص.\ \pageref{sub:communication-model}
  \item[نایکوئیست و هارتلی] \hyperref[sub:nyquist-and-hartley]{پیش از شانون: نایکوئیست و هارتلی}، ص.\ \pageref{sub:nyquist-and-hartley}
\end{description}
